% titel: Blatt 1, A6: Knobelaufgabe: Nullen von 100!

% Teil I: Werkzeuge, kein Lösungsweg.

% Grundidee: Eine Endnull entsteht aus 2*5=10. In n! gibt es viel mehr
% Zweien als Fünfen, also zählt man die Fünfen (Primfaktor 5).

Warum hat $10!=3628800$ genau zwei Endnullen? Eine Null am Ende ist
das Ergebnis der Multiplikation mit $10=2\cdot5$. In einer Fakultät
stecken beide Primfaktoren in Hülle und Fülle, aber \emph{nie gleich
viel}: Es gibt deutlich mehr Zweien als Fünfen. Deshalb ist die Zahl der
Fünfen der Engpass --- jede Fünf ``paart`` sich mit einer der vielen
Zweien und erzeugt eine $10$, also eine Endnull.

\begin{werkzeugkasten}[Endnullen einer Fakultät zählen]
Die Anzahl der Primfaktoren $p$ in $n!$ zählt die sog. Legendre-Formel:
\begin{equation}
  \left\lfloor\frac{n}{5}\right\rfloor+
  \left\lfloor\frac{n}{25}\right\rfloor+
  \left\lfloor\frac{n}{125}\right\rfloor+\dots
\end{equation}
(alle $5^k\le n$). \textbf{Warum die Potenzen?} Jede Zahl, die ein
Vielfaches von $25$ ist, enthält eine \emph{zusätzliche} Fünf, die von
$\lfloor n/5\rfloor$ nur einmal erfasst wird --- deshalb zählt man ab
$25$ noch einmal nach, usw. Für $n=10$: $\lfloor10/5\rfloor=2$,
$\lfloor10/25\rfloor=0$, Summe $2$ --- genau die zwei Endnullen von
$10!$.
\end{werkzeugkasten}

\begin{achtung}[Typische Fehler]
\begin{itemize}
  \item \textbf{Nicht die Zweien, sondern die Fünfen zählen} --- Nullen
        entstehen durch Fünf in Kombination mit genügend Zweien; da die
        Zweien immer im \"Uberschuss sind, ist die Fünf der
        limitierende Faktor.
  \item Jede Potenz $25,125,\dots$ gibt einen \emph{weiteren}
        Primfaktor $5$ und wird mitgezählt.
  \item Beim konkreten Z\"ahlen nicht vergessen, dass z.\,B. bei $100!$
        die Zahlen $25,50,75,100$ je eine extra Fünf tragen (je zwei
        f\"unf-teilbare Stufen).
\end{itemize}
\end{achtung}
